\scanpage{27}
\mspage{21}

Prouvons $3^\circ$. Supposons d'abord qu'il existe
un quadruplet ``clos'' $(C,TC,F,TF)$, s'insérant
dans un ``pré-quintuplet de Quillen clos'' qui
soit une \inserted{$Q$-}cat. de modèles à droite, i.e.
\[
 TF=F\cap W=F\cap W_{\mathrm{univ}}.
\]

\deleted{Comme} Prouvons donc
\[
 \deleted{F\subset\mathrm{Fib}_W,\ \text{on a}\quad TF=}
\]
\deleted{\illegible{3}}
\[
 \text{que}\qquad\underbrace{(W^{\mathrm{univ}})_*\cap W_{\mathrm{univ}}}_{A_W}
 \quad\text{\deleted{\unsure{est clos}} est clos à \unsure{droite}}
\]

En effet, on a
\[
 TC\subset W^{\mathrm{univ}}
\]
d'où
\[
 (W^{\mathrm{univ}})_*\subset(TC)_*=F
\]
donc
\[
 A=(W^{\mathrm{univ}})_*\cap W_{\mathrm{univ}}\subset F\cap W_{\mathrm{univ}}
 \underset{\substack{\uparrow\\\text{(par hyp.)}}}{=}TF,
\]
or $TF$ est clos à droite et contenu dans
$W_{\mathrm{univ}}$, d'où
\[
 \overline A\subset\deleted{\overline{\unsure{(W^{\mathrm{univ}})_*\cap W_{\mathrm{univ}}}}}\subset TF\subset W_{\mathrm{univ}},
\]
et comme $(W^{\mathrm{univ}})_*$ est clos à droite, on a aussi
\[
 \deleted{\overline{\unsure{(W^{\mathrm{univ}})_*\cap W_{\mathrm{univ}}}}}\overline A\subset(W^{\mathrm{univ}})_*,
\]
\deleted{q.e.d.}
% slanted
\inserted{d'où $\overline A\subset W_{\mathrm{univ}}\cap(W^{\mathrm{univ}})_*=A$, donc $A$ est clos à droite, qed.}

Prouvons la réciproque, i.e. supposons
satisfaites les conditions du $2^\circ$, \uline{plus} la
condition $(*)$ ci-dessus. On doit construire
\unsure{le|les} \unsure{qu. demandé|quadruplet}.

\deleted{On prendra
\[
 \begin{aligned}
 TC&=W^{\mathrm{univ}}&&\deleted{\unsure{TC}}\ \text{(le plus grand possible)}\\
 C&=\deleted{\unsure{\mathrm{Cof}_W}}\widetilde{(W_{\mathrm{univ}})^*\cup W^{\mathrm{univ}}}
 &&\text{(partie close}\\
 &&&\text{à g. engendrée par }(W_{\mathrm{univ}})^*\text{ et }W^{\mathrm{univ}}\text{)}
 \end{aligned}
\]
d'où $F=(TC)_*=(W^{\mathrm{univ}})_*$}
% archive sheet number: 26
